\documentclass[10pt]{article}
\usepackage[T1]{fontenc}
\usepackage[utf8]{inputenc}
\usepackage[french]{babel}
\usepackage{amsmath,amssymb,mathrsfs,mathpazo}
\usepackage{geometry}
\geometry{a4paper,margin=2cm}
\usepackage{graphicx,array,tabularx,tabularray,colortbl,multirow,multicol,booktabs,enumitem,pifont,ulem,fancybox,diagbox,cancel}
\usepackage{pstricks,pstricks-add,pst-plot,pst-tree,pst-node,pst-text,pst-func,pst-math,pst-bspline,pst-3dplot,pst-eucl,pst-solides3d}
\let\pscalculate\undefined
\usepackage{pst-bezier}
\usepackage{tikz,pgfplots,tkz-tab}
\usetikzlibrary{arrows,arrows.meta,patterns,calc,positioning,trees,decorations.markings,shapes.geometric,shapes.misc,shapes.arrows,intersections,decorations.pathmorphing,decorations.pathreplacing,backgrounds,shadows,babel}
\pgfplotsset{compat=1.18}
\usepgfplotslibrary{fillbetween}
\usepackage[np]{numprint}
\usepackage{listings,makecell,pas-tableur,twoopt}
\let\barre\undefined
\usepackage[pstricks]{tablvar}
\definecolor{darkgreen}{RGB}{0,100,0}
\providecommand{\red}{\color{red}}
\providecommand{\blue}{\color{blue}}
\providecommand{\green}{\color{green}}
\providecommand{\black}{\color{black}}
\providecommand{\white}{\color{white}}
\providecommand{\gray}{\color{gray}}
\providecommand{\cyan}{\color{cyan}}
\providecommand{\magenta}{\color{magenta}}
\providecommand{\yellow}{\color{yellow}}
\providecommand{\violet}{\color{violet}}
\providecommand{\orange}{\color{orange}}
\providecommand{\euro}{\texteuro}
\providecommand{\EUR}{\texteuro}
\providecommand{\ieme}{\textsuperscript{e}}
\providecommand{\degres}{\ensuremath{^\circ}}
\providecommand{\degree}{\ensuremath{^\circ}}
\providecommand{\vv}[1]{\overrightarrow{#1}}
\providecommand{\Coord}[1]{#1}
\providecommand{\decosix}{}
\providecommand{\starredbullet}{\ensuremath{\bullet}}
\pagestyle{empty}
\setlength{\parindent}{0pt}
\usepackage{adjustbox}\DeclareUnicodeCharacter{2236}{:}\DeclareUnicodeCharacter{2190}{\ensuremath{\leftarrow}}
\providecommand{\R}{\mathbb{R}}
\providecommand{\N}{\mathbb{N}}
\providecommand{\Z}{\mathbb{Z}}
\providecommand{\Q}{\mathbb{Q}}
\providecommand{\D}{\mathbb{D}}
\providecommand{\C}{\mathbb{C}}
\providecommand{\e}{{\mathrm{e}}}
\providecommand{\vect}[1]{\overrightarrow{#1}}
\providecommand{\vectt}[1]{\overrightarrow{\mathrm{#1}}}
\providecommand{\vvt}[1]{\overrightarrow{\mathrm{#1}}}
\providecommand{\np}[2][]{{#2~#1}}
\providecommand{\vv}[1]{{\overrightarrow{#1}}}
\renewcommand{\d}{{\mathrm{d}}}
\providecommand{\degres}{{^\circ}}
\providecommand{\degre}{{^\circ}}
\providecommand{\blue}{\color{blue}}
\providecommand{\red}{\color{red}}
\providecommand{\enskip}{\quad}
\providecommand{\llbracket}{\lbrack\!\lbrack}
\providecommand{\rrbracket}{\rbrack\!\rbrack}
\providecommand{\euro}{{\text{€}}}
\providecommand{\no}{{\text{n°}}}
\providecommand{\interval}[2]{{[#1;#2]}}
\providecommand{\intervFO}[2]{{[#1;#2[}}
\providecommand{\dotfill}{\ldots}

\providecommand{\ds}{\displaystyle}
\providecommand{\textding}[1]{#1}
\providecommand{\Oij}{\(\mathrm{O};\vec{\imath},\vec{\jmath}\)}
\providecommand{\Oijk}{\(\mathrm{O};\vec{\imath},\vec{\jmath},\vec{k}\)}
\providecommand{\Ouv}{\(\mathrm{O};\vec{u},\vec{v}\)}
\providecommand{\intervFF}[2]{\ensuremath{[#1;#2]}}
\renewcommand{\euro}{\ifmmode\text{\texteuro}\else\texteuro\fi}
\DeclareUnicodeCharacter{00B0}{\ifmmode^{\circ}\else\textdegree\fi}
\DeclareUnicodeCharacter{20AC}{\ifmmode\text{\texteuro}\else\texteuro\fi}
% Énoncés : Éducation nationale. Transcription et corrigés : APMEP, Denis Vergès et contributeurs. Découpage et classement : CoopMaths, Sébastien Lozano, Éric Elter et contributeurs. Adaptation de la présentation : Maths Explorer.
% CC BY-SA : https://coopmaths.fr/www/about/
% https://coopmaths.fr/alea/static/eam/2025/tex/eam_specifique_2025_06_sujet0v2_3_cor.tex
\begin{document}
\noindent\textit{APMEP / CoopMaths. CC BY-SA. Présentation adaptée par Maths Explorer.}\par\medskip
\medskip

\textbf{Partie A}

\medskip

\begin{enumerate}
\item On complète le tableau des effectifs :

\begin{center}
\begin{tabularx}{0.5\linewidth}{|l|*{3}{>{\centering \arraybackslash}X|}}\cline{2-4}
\multicolumn{1}{c|}{~} & Fille & Garçon&Total \\ \hline
Anglais & 712 & 728&\np{1440} \\ \hline
Autre LV1 & 288 & 272&\np{560} \\ \hline
Total&\np{1000}&\np{1000}&\np{2000} \\ \hline
\end{tabularx}
\end{center}

Il y a donc autant de filles que de garçons.
\item 712 filles ont choisi l'Anglais, donc $p(A \cap F) = \dfrac{712}{\np{2000}}$.
\item On a $p_F(A) = \dfrac{p(F \cap A)}{p(F)} = \dfrac{p(A \cap F)}{p(F)} = \dfrac{\dfrac{712}{\np{2000}}}{\dfrac{1000}{\np{2000}}} = \dfrac{712}{\np{1000}}$.
\item On a $p(A) \times p(F) = \dfrac{712}{\np{2000}} \times \dfrac12 = \dfrac{712}{\np{4000}} = \dfrac{356}{\np{2000}}$.

Donc $p(A) \times p(F) \ne p(A \cap F)$ : les évènements $A$ et $F$ ne sont pas indépendants.
\item 728 élèves garçons ont choisi l'anglais et 272 ne l'ont pas choisi : or $728 \ne 3 \times 272$ car ce produit se termine par un 6.

L'affirmation est fausse.
\end{enumerate}

\bigskip

\textbf{Partie B}

\medskip

\begin{enumerate}
\item On a $p(P) = \dfrac{1}{4}$, donc $p(F) = 1 - p(P) = 1 - \dfrac{1}{4} = \dfrac34$.
\item~

\begin{center}
\pstree[treemode=R,nodesep=1ex,levelsep=2.5cm]{\TR{}}
{\pstree{\TR{$P~~$}\taput{\small $\frac14$}}
		{\pstree{\TR{$P~~$}\taput{\small $\frac14$}}
			{\TR{$P~~$}\taput{\small $\frac14$}
			\TR{$F~~$}\tbput{\small $\frac34$}
			}
		\pstree{\TR{$F~~$}\tbput{\small $\frac34$}}
			{\TR{$P~~$}\taput{\small $\frac14$}
			\TR{$F~~$}\tbput{\small $\frac34$}
			}
		}
\pstree{\TR{$F~~$}\tbput{\small $\frac34$}}
		{\pstree{\TR{$P~~$}\taput{\small $\frac14$}}
			{\TR{$P~~$}\taput{\small $\frac14$}
			\TR{$F~~$}\tbput{\small $\frac34$}
			}
		\pstree{\TR{$F~~$}\tbput{\small $\frac34$}}
			{\TR{$P~~$}\taput{\small $\frac14$}
			\TR{$F~~$}\tbput{\small $\frac34$}
			}
	}
}
\end{center}

\item La probabilité d'obtenir une fois Pile est :

$p(PFF) + p(FPF) + p(FFP) = \frac14 \times \frac34 \times \frac34 + \frac34 \times \frac14 \times \frac34 + \frac34 \times \frac34 \times \frac14 = \frac{9}{64} + \frac{9}{64} + \frac{9}{64} = \dfrac{27}{64}$.
\item La probabilité de ne jamais obtenir pile est égale à :

$p(FFF) = \dfrac34 \times \dfrac34 \times \dfrac34 = \dfrac{27}{64}$ (branche du bas).
\end{enumerate}

\end{document}
